\section*{الفصل \ref{ch:b2:euler-bernoulli} --- الموائع المثالية: أويلر وبرنولي}

\begin{solution}{exo:b2:euler-bernoulli:1}
$v_2 = 1.5 \times 4 = \qty{6.0}{m/s}$، و $P_2 = 3.0 \times 10^5 - 500(36 - 2.25) =
\qty{2.83}{bar}$. وعند \qty{2}{cm}: $v = \qty{24}{m/s}$، و $P = 3.0 \times 10^5 - 500
(576 - 2.25) = \qty{0.13}{bar}$ --- وهو قريب من ضغط البخار. والنتيجة
السالبة تعني أن التدفق المفترض مستحيل: فيتكهّف الماء
ويُخنق الجريان.
\end{solution}

\begin{solution}{exo:b2:euler-bernoulli:2}
$v = \sqrt{2 \times 15000/0.74} = \qty{201}{m/s}$؛ والمؤشَّرة $\sqrt{2 \times 15000/
1.22} = \qty{157}{m/s}$. والسرعة الجوية المؤشَّرة تقيس \omterm{thm:b2:euler-bernoulli:bernoulli}{الضغط الديناميكي}،
وهو ما يحسّه الجناح: فسرعة الانهيار سرعة \emph{مؤشَّرة} ثابتة
على أي ارتفاع.
\end{solution}

\begin{solution}{exo:b2:euler-bernoulli:3}
على عمق \qty{2.0}{m}: $v = \sqrt{2g \times 2} = \qty{6.3}{m/s}$؛ و $D_V = 0.6 \times 10^{-4}
\times 6.3 = \qty{0.38}{L/s}$؛ وزمن السقوط $\sqrt{2/g} = \qty{0.45}{s}$، والمدى
\qty{2.8}{m}. والمدى $= 2\sqrt{h(H - h)}$ متناظر في $h \leftrightarrow H -
h$: فثقب على عمق \qty{1.0}{m} (أي \qty{2}{m} فوق الأرض) يصيب النقطة
نفسها.
\end{solution}

\begin{solution}{exo:b2:euler-bernoulli:4}
$\Delta P = \tfrac12 \times 1.2 \times 900 = \qty{540}{Pa}$، والضغط في الخارج أدنى:
فالقوة الصافية \qty{54}{kN} \emph{نحو الأعلى}؛ ويزن السقف \qty{49}{kN}: فيُقتلع.
(والسقوف الحقيقية تنهار من حوافها أولًا، حيث ينفصل الجريان.)
\end{solution}

\begin{solution}{exo:b2:euler-bernoulli:5}
(أ) $\tan\theta = 3/9.81$، و $\theta = \qty{17}{\degree}$؛ و $2 \times 0.306 = \qty{0.61}{m}$
(أعلى في الخلف). (ب) $\Delta P = \rho aL = \qty{6.0}{kPa}$. (ج) $\tan\theta =
0.61$، والفرق \qty{1.22}{m}: فيرتفع الأمام \qty{0.61}{m} فوق المتوسط
\qty{0.40}{m}: أي \qty{1.01}{m} $>$ \qty{0.8}{m}، فيصطدم بالغطاء. (د) $a = v^2/R =
\qty{5.6}{m/s^2}$: ومنه $\theta = \qty{30}{\degree}$، والجانب الخارجي أعلى.
\end{solution}

\begin{solution}{exo:b2:euler-bernoulli:6}
(أ) مجسّم مكافئ $z = z_0 + \Omega^2r^2/2g$، مع $\Omega = 4\pi = \qty{12.6}{rad/s}$. (ب)
$\Omega^2R^2/2g = 158 \times 0.0225/19.6 = \qty{0.18}{m}$. (ج) ويقع الارتفاع الوسطي
للمجسّم المكافئ فوق القرص في المنتصف: فالمركز $0.20 - 0.09 = \qty{0.11}{m}$، والحافة
\qty{0.29}{m}. (د) ويجفّ المركز حين $\Omega^2R^2/4g = h_0$: أي $\Omega = \sqrt{4gh_0}/R =
\qty{18.7}{rad/s} = 3.0$ دورات في الثانية. (ه) والمجسّم المكافئ يجمع
الأشعة المتوازية في نقطة واحدة، عند $f = g/2\Omega^2$؛ ومقاريب المرآة
السائلة تُدير الزئبق.
\end{solution}

\begin{solution}{exo:b2:euler-bernoulli:7}
(أ) $v = \sqrt{2g \times 3} = \qty{7.7}{m/s}$؛ و $D_V = \pi \times 10^{-4} \times 7.7 =
\qty{2.4}{L/s}$. (ب) $P = P_0 - \rho g(2) - \tfrac12\rho v^2 = 100 - 19.6 - 29.4 =
\qty{51}{kPa}$. (ج) $P_0 - \rho gh_c - \tfrac12\rho v^2 \ge \qty{2.3}{kPa}$: أي $h_c \le
\qty{7.0}{m}$. (د) والهواء انضغاطي وخفيف: فلا عمود سائل متصل ينقل
الضغط، ويبقى الماء على الجانبين حيث هو.
\end{solution}

\begin{solution}{exo:b2:euler-bernoulli:8}
(أ) $\Delta P = mg/S = \qty{654}{Pa}$؛ و $\tfrac12\rho v^2 = \qty{2.2}{kPa}$، و $C_L = 0.30$.
(ب) $v_{\text{فوق}}^2 = 58^2 + 2 \times 654/1.22$: أي \qty{66.6}{m/s}. (ج) $F_L' =
\qty{981}{N/m} = \rho v\Gamma$: ومنه $\Gamma = \qty{13.4}{m^2/s}$. (د) $v = \sqrt{2mg/
\rho SC_L} = \sqrt{19620/25.6} = \qty{28}{m/s}$.
\end{solution}

\begin{solution}{exo:b2:euler-bernoulli:9}
(أ) $v_2 = \qty{2.5}{m/s}$؛ و $\Delta P = \tfrac12 \times 1060 \times (6.25 - 0.25) =
\qty{3.2}{kPa}$، و $P_2 = \qty{9.8}{kPa}$. (ب) وهو ما يزال فوق \qty{1}{kPa}: فلا. (ج) وعند
\qty{0.05}{cm^2} يكون $v = \qty{5}{m/s}$ و $\Delta P = \qty{13}{kPa}$: أي $P_2 \approx 0$، دون
ضغط النسيج --- فينهار الجدار، ويتوقف الجريان، ويستعيد
الضغط قيمته، فينفتح من جديد: وهي رفرفة (وهي اللغط الذي يسمعه الطبيب).
\end{solution}

\begin{solution}{exo:b2:euler-bernoulli:10}
(أ) $c_PT_0 = c_PT + v^2/2$ مع $c_P = \gamma R/(\gamma - 1)M_{\text{مول}}$ و
$c^2 = \gamma RT/M_{\text{مول}}$: ومنه $v^2/2c_PT = \tfrac12(\gamma - 1)M^2$. (ب) عند $M =
0.85$: $T_0 = 220 \times 1.145 = \qty{252}{K}$؛ وعند $M = 2$: $220 \times 1.8 = \qty{396}{K}$؛
وعند $M = 25$: $220 \times 126 = \qty{28000}{K}$ --- وهو بلا معنى: إذ يتفكك الهواء
ويتأيّن، ولا يبقى $c_P$ ثابتًا (ودرجات حرارة الركود الحقيقية
بضعة آلاف كلفن). (ج) $P_0/P = (1 + \tfrac{\gamma - 1}2M^2)^{\gamma/(\gamma
- 1)} \approx 1 + \tfrac{\gamma}2M^2 + \tfrac{\gamma}8M^4$، و $\gamma PM^2 = \rho v^2$:
ومنه $P_0 - P = \tfrac12\rho v^2(1 + M^2/4)$. وقراءة $v$ من $\sqrt{2(P_0 - P)/\rho}$
تفرط في تقديرها بالنسبة $\sqrt{1 + M^2/4} - 1 \approx 9\%$ عند $M = 0.85$.
\end{solution}

\begin{solution}{exo:b2:euler-bernoulli:11}
(أ) من السطح ($v \approx 0$، و $P_0$، وارتفاع $h$) إلى المخرج ($v$، و $P_0$،
وارتفاع $0$)، مع $\int\partial_tv\,\dd s = \ell\,\dd v/\dd t$. (ب) $\dd v/(V^2 -
v^2) = \dd t/2\ell$، و $V^2 = 2gh$: ومنه $\operatorname{artanh}(v/V) = Vt/2\ell$. (ج) $V =
\qty{9.9}{m/s}$؛ و $\tanh x = 0.9$ عند $x = 1.47$: أي $t = 2\ell x/V = \qty{15}{s}$. (د)
$\dd v/\dd t \approx \qty{-20}{m/s^2}$: ومنه $\Delta P = \rho\ell \times 20 = \qty{9.9}{bar}$.
والماء انضغاطي قليلًا: فالإغلاق السريع يطلق موجة ضغط،
$\Delta P = \rho c\Delta v \approx \qty{140}{bar}$.
\end{solution}

\begin{solution}{exo:b2:euler-bernoulli:12}
(أ) مستقر ومثالي ولاانضغاطي ولادوراني: فثابت واحد
للجريان كله. (ب) وعلى السطح $P = P_0$: أي $\tfrac12\rho v^2 + \rho gz =
0$، ومنه $z = -\Gamma^2/8\pi^2gr^2$. (ج) \omterm{thm:b2:euler-bernoulli:euler}{ومعادلة أويلر} قطريًّا: $\partial_rP = \rho\Omega^2r$،
وهيدروستاتيكيًّا شاقوليًّا: فيحقق السطح $z(r) - z(a) = \Omega^2(r^2 -
a^2)/2g$. (د) الحوض: $\Omega = \Gamma/2\pi a^2 = \qty{191}{rad/s}$؛ و $z(a) = -\Omega^2a^2/
2g = \qty{-4.7}{cm}$، ويضيف اللبّ \qty{4.7}{cm} أخرى: فالعمق \qty{9.3}{cm}.
والإعصار القمعي: $\Delta P = \tfrac12\rho v_{\max}^2$ في الخارج ومثلها في الداخل:
أي $\rho v_{\max}^2 = 1.2 \times 6400 = \qty{7.7}{kPa}$ دون الضغط المحيط عند المركز.
\end{solution}

\begin{solution}{pb:b2:euler-bernoulli:1}
\textbf{1.} $\rho g \times 150 = \qty{14.7}{bar}$ فوق الضغط الجوي؛ و $\tfrac12\rho gH_d^2
\times \qty{1}{m} = \qty{1.1e8}{N}$.

\textbf{2.} $\rho gH = \qty{19.6}{bar}$ نسبيًّا، و \qty{20.6}{bar} مطلقًا.

\textbf{3.} $v = 60/7.07 = \qty{8.5}{m/s}$؛ و $\tfrac12\rho v^2 = \qty{0.36}{bar}$؛ و $P =
P_0 + \rho gH - \tfrac12\rho v^2 = \qty{20.3}{bar}$ مطلقًا.

\textbf{4.} $v_j = \sqrt{2gH} = \qty{62.6}{m/s}$؛ و $s = 60/62.6 = \qty{0.96}{m^2}$، و $d =
\qty{1.1}{m}$.

\textbf{5.} $\tfrac12\rho D_Vv_j^2 = \tfrac12\rho D_V(2gH) = \rho gD_VH = \qty{118}{MW}$.

\textbf{6.} $m = \rho S\ell = \qty{2.8e6}{kg}$، و $E_k = \tfrac12mv^2 = \qty{1.0e8}{J}$؛
وزمن العبور $\ell/v = \qty{47}{s}$.

\textbf{7.} $\rho gH = (P - P_0) + \tfrac12\rho v^2$: أي \qty{19.3}{bar} من الضغط
و \qty{0.36}{bar} من الطاقة الحركية؛ وفي الفوهة يهبط الضغط
إلى $P_0$ بينما يرتفع $\tfrac12\rho v^2$ إلى $\rho gH$: فبُعيدها $P = P_0$.

\textbf{8.} $\Delta h = \lambda(\ell/D)v^2/2g = 0.015 \times 133 \times 3.67 = \qty{7.3}{m}$،
أي $3.7\%$ من الاستطاعة.

\textbf{9.} $v = \qty{19.1}{m/s}$، و $v^2/2g = \qty{18.6}{m}$، و $\Delta h = 0.015 \times 200 \times
18.6 = \qty{56}{m}$: أي $28\%$ ضائعة --- فالتوفير في الفولاذ يُدفع في كل
ثانية طاقةً.

\textbf{10.} $v_t = 60/\pi = \qty{19.1}{m/s}$؛ و $\Delta P = 500(19.1^2 - 8.49^2) =
\qty{146}{kPa}$.

\textbf{11.} $\Delta P = (\rho_{\text{زئبق}} - \rho)g\,\Delta h$: ومنه $\Delta h = 146000/(12600
\times 9.81) = \qty{1.18}{m}$.

\textbf{12.} $\Delta P \propto v^2 \propto D_V^2$: ومنه $D_V \propto \sqrt{\Delta h}$؛ ونصف
التدفق يعطي الربع، أي \qty{0.30}{m}.

\textbf{13.} المأخذ البارز يواجه الجريان أو يزعجه فيقرأ جزءًا من
\omterm{thm:b2:euler-bernoulli:bernoulli}{الضغط الديناميكي} (\omterm{prop:b2:euler-bernoulli:pitot}{كأنبوب بيتو})، لا الضغط السكوني.

\textbf{14.} $P = P_0 - \rho g \times 6 - \tfrac12\rho v^2 = 100 - 59 - 36 = \qty{5}{kPa}$:
أي قرب ضغط البخار، فيخرج الهواء المنحلّ وقد يغلي الماء
(تكهّف)، فينكسر العمود. فليُخفَض ارتفاع القمة، أو يُركَّب صمام
هواء ومضخة تخلية، أو يُخفَّض التدفق.

\textbf{15.} $D_V/A = 60/2 \times 10^6 = \qty{3e-5}{m/s} = \qty{11}{cm/h}$.

\textbf{16.} $A\,\dd h/\dd t = -s\sqrt{2gh}$: ومنه $\sqrt h = \sqrt{h_0} - (s/A)\sqrt{g/2}\,t$.

\textbf{17.} $\Delta t = (A/s)\sqrt{2/g}(\sqrt{200} - \sqrt{190}) = 2.09 \times 10^6 \times
0.45 \times 0.36 = \qty{3.4e5}{s} = \qty{3.9}{d}$؛ وعند التدفق التصميمي الثابت
$10A/D_V = \qty{3.3e5}{s}$: أي القيمة نفسها تقريبًا، إذ عوّلت الفوهة
على جريان توريتشيلي عند $h \approx H$.

\textbf{18.} $\ell\,\dd v/\dd t = gH - v^2/2$: ومنه $v = V\tanh(Vt/2\ell)$، و $V = \qty{62.6}{m/s}$،
وثابت الزمن $2\ell/V = \qty{13}{s}$.

\textbf{19.} عند $v = V/2$: $\dd v/\dd t = (gH - V^2/8)/\ell = \qty{3.7}{m/s^2}$؛
و $P = P_0 + \rho gH - \tfrac12\rho v^2 - \rho\ell\,\dd v/\dd t = 1.0 + 19.6 - 4.9 - 14.7
= \qty{1.0}{bar}$: أي إن معظم العلو منشغل بتسريع العمود، ومنه
يكون الضغط في الأسفل أدنى بكثير من قيمته المستقرة.

\textbf{20.} $|\dd v/\dd t| = 8.49/10 = \qty{0.85}{m/s^2}$: ومنه $\Delta P = \rho\ell \times 0.85 =
\qty{3.4}{bar}$.

\textbf{21.} $\Delta P = 1000 \times 1400 \times 8.49 = \qty{119}{bar}$، أي ستة أضعاف
الضغط السكوني. وذهاب الموجة وإيابها $2\ell/c = \qty{0.57}{s}$: فالإغلاق
الأسرع من ذلك «مفاجئ» وينال فرط ضغط جوكوفسكي كاملًا؛
وأما الأبطأ فينال أقلّ.

\textbf{22.} $Av = A_s\,\dd z/\dd t$.

\textbf{23.} أويلر على امتداد النفق، من الخزان (وضغطه
$P_0 + \rho g\times$ العمق) إلى قدم البئر (وضغطه $P_0 + \rho g
(\text{العمق} + z)$): $\rho L\,\dd v/\dd t = -\rho gz$؛ ومع $v = (A_s/A)\dot z$:
$\ddot z + (gA/LA_s)z = 0$؛ ومنه الدور $T = 2\pi\sqrt{LA_s/gA}$، أي $2\pi\sqrt{2000 \times 100/
68.7} = \qty{340}{s} \approx \qty{5.7}{min}$.

\textbf{24.} $\tfrac12\rho ALv^2 = \tfrac12\rho gA_sz_{\max}^2$ مع $v = 60/7 =
\qty{8.6}{m/s}$: ومنه $z_{\max} = v\sqrt{AL/gA_s} = \qty{32}{m}$.

\textbf{25.} السدّ: هيدروستاتيك، $\rho gh$. والعنفة وصماماتها
مغلقة: الأمر نفسه، $\rho gH$. وأنبوب الجرّ في التشغيل: مبرهنة
برنولي، $\tfrac12\rho v^2$ مأخوذة من العلو السكوني. والصدمة
المائية: الانضغاطية، $\rho c\Delta v$. وارتفاع البئر: \omterm{thm:b2:euler-bernoulli:euler}{معادلة أويلر}
غير المستقرة على امتداد النفق، واهتزاز دوره $2\pi\sqrt{LA_s/gA}$.
\end{solution}
